\chapter{TEORI FREDHOLM}

\section{Alternatif Fredholm}

Pada tahun 1903, Erik Ivar Fredholm menerbitkan sebuah makalah rintisan tentang persamaan integral dalam jurnal
\emph{Acta Mathematica}. Di antara banyak hasil pentingnya terdapat teorema yang kini kita kenal sebagai
\emph{alternatif Fredholm}. Kita menyatakan suatu versi hasil ini dalam bahasa
yang tersedia bagi Fredholm pada masa itu.

\begin{prop}[Alternatif Fredholm I]\label{004011} Tetapkan $\lambda \in \C\setminus\{0\}$ dan misalkan $k$ suatu fungsi kontinu bernilai kompleks
pada persegi satuan $[0,1] \times [0,1]$. \textbf{Entah} persamaan-persamaan
 \index{alternatif Fredholm!versi I}%
tak homogen
 \begin{align}
   \lambda f(s) &- \int_0^1 k(s,t)f(t)\,dt = g(s)\label{004011i} \text{ dan}    \tag{1}\\
   \conj{\lambda} h(s) &- \int_0^1 \conj{k(t,s)}h(t)\,dt = j(s)\label{004011ii} \tag{2}
 \end{align}
memiliki solusi $f$ dan $h$ untuk setiap $g$ dan $j$ yang diberikan, masing-masing, dengan solusi-solusi tersebut
tunggal; dalam hal ini persamaan homogen yang bersesuaian
 \begin{align} \lambda f(s) &- \int_0^1 k(s,t)f(t)\,dt = 0\label{004011iii} \text{ dan} \tag{3}\\
        \conj{\lambda} h(s) &- \int_0^1 \conj{k(t,s)}h(t)\,dt = 0\label{004011iv} \tag{4}
 \end{align}
hanya memiliki solusi trivial; \textbf{---atau---}\\
persamaan homogen \eqref{004011iii} dan \eqref{004011iv} masing-masing memiliki
sejumlah hingga (dan tak nol) solusi bebas linear yang sama banyaknya, yaitu $f_1,\dots,f_n$ dan $h_1,\dots,h_n$,
secara berturut-turut; dalam hal ini persamaan tak homogen \eqref{004011i} dan
\eqref{004011ii} memiliki solusi jika dan hanya jika $g$ dan $j$ memenuhi
 \begin{align} \int_0^1 h_k(t)\conj{g(t)}\,dt &= 0\label{004011v} \text{ dan} \tag{5}\\
              \int_0^1 j(t) \conj{f_k(t)}\,dt &= 0\label{004011vi} \tag{6}
 \end{align}
untuk $k = 1, \dots, n$.
\end{prop}

Pada tahun 1906, David Hilbert telah menyadari bahwa pengintegralan itu sendiri sangat sedikit berkaitan dengan
kebenaran hasil ini. Ia menemukan bahwa yang penting adalah kekompakan
operator integral yang dihasilkan. (Pada awal abad ke-$20$
\emph{kekompakan} disebut \emph{kontinuitas lengkap}. Istilah \emph{ruang Hilbert}
telah digunakan sejak tahun 1911 sehubungan dengan ruang barisan~$l_2$. Baru pada
tahun 1929 John von Neumann memperkenalkan gagasan---dan aksioma yang
mendefinisikan---\emph{ruang Hilbert abstrak}.) Jadi, berikut suatu versi yang agak lebih mutakhir dari
\emph{alternatif Fredholm}.

\begin{prop}[Alternatif Fredholm II]\label{004031} Jika $K$ suatu operator kompak pada ruang
 \index{alternatif Fredholm!versi II}%
Hilbert, $\lambda \in \C\setminus\{0\}$, dan $T = \lambda I - K$, maka \textbf{entah} persamaan-persamaan
tak homogen
 \begin{align}
            Tf &= g\label{004031i} \text{ dan} \tag{1'}\\
            T^* h &= j\label{004031ii} \tag{2'}
 \end{align}
memiliki solusi $f$ dan $h$ untuk setiap $g$ dan $j$ yang diberikan, masing-masing, dengan solusi-solusi tersebut
tunggal; dalam hal ini persamaan homogen yang bersesuaian
 \begin{align}
            Tf &= 0\label{004031iii} \text{ dan} \tag{3'}\\
          T^*h &= 0\label{004031iv} \tag{4'}
 \end{align}
hanya memiliki solusi trivial; \textbf{---atau---}\\
persamaan homogen \eqref{004031iii} dan \eqref{004031iv} masing-masing memiliki
sejumlah hingga (dan tak nol) solusi bebas linear yang sama banyaknya, yaitu $f_1,\dots,f_n$ dan $h_1,\dots,h_n$,
secara berturut-turut; dalam hal ini persamaan tak homogen \eqref{004031i} dan
\eqref{004031ii} memiliki solusi jika dan hanya jika $g$ dan $j$ memenuhi
 \begin{align} h_k &\perp g\label{004031v} \text{ dan} \tag{5'}\\
                 j &\perp f_k\label{004031vi} \tag{6'}
 \end{align}
untuk $k = 1, \dots, n$.
\end{prop}

Perhatikan bahwa dengan menggunakan beberapa fakta elementer mengenai kernel dan jangkauan
operator serta keortogonalan dalam ruang Hilbert, kita dapat meringkas pernyataan
dari~\ref{004031} secara cukup banyak.

\begin{prop}[Alternatif Fredholm IIIa]\label{004033} Jika $T = \lambda I - K$ dengan
 \index{alternatif Fredholm!versi IIIa}%
$K$ suatu operator kompak pada ruang Hilbert dan $\lambda \in \C\setminus\{0\}$, maka
 \begin{enumerate}
  \item $T$ injektif jika dan hanya jika bersifat surjektif,
  \item $\ran T^* = (\ker T)^\perp$, dan
  \item $\dim\ker T = \dim \ker T^*$.
 \end{enumerate}
Selain itu, syarat (1) dan (2) berlaku untuk $T^*$ maupun~$T$.
\end{prop}


















\section{Alternatif Fredholm -- lanjutan}
\begin{defn}\label{004034} Suatu operator $T$ pada ruang Banach disebut
 \index{operator Riesz--Schauder}%
 \index{operator!Riesz--Schauder}%
\df{operator Riesz--Schauder} jika dapat ditulis dalam bentuk $T = S + K$ dengan $S$
invertibel dan $K$ kompak.
\end{defn}

Materi dalam bagian~\ref{sec_onbases} dan definisi sebelumnya memungkinkan kita
memperumum versi \emph{alternatif Fredholm} yang diberikan dalam~\ref{004033} ke ruang
Banach.

\begin{prop}[Alternatif Fredholm IIIb]\label{0040343} Jika $T$ suatu operator Riesz--Schauder
 \index{alternatif Fredholm!versi IIIb}%
pada ruang Banach, maka
 \begin{enumerate}
  \item $T$ injektif jika dan hanya jika bersifat surjektif,
  \item $\ran T^* = (\ker T)^\perp$, dan
  \item $\dim\ker T = \dim \ker T^* < \infty$.
 \end{enumerate}
Selain itu, syarat (1) dan (2) berlaku untuk $T^*$ maupun~$T$.
\end{prop}

\begin{prop} Jika $M$ subruang tertutup dari ruang Banach $B$, maka $M^\perp \cong (B/M)^*$.
\end{prop}

\begin{proof} Lihat \cite{Conway:1990}, halaman 89. \ns \end{proof}

\begin{defn} Jika $T\colon V \sto W$ suatu pemetaan linear antara ruang vektor, maka
 \index{kokernel}%
\df{kokernel} pemetaan tersebut didefinisikan oleh
  \[ \coker T = W/ \ran T\,. \]
Ingat bahwa
 \index{kodimensi}%
\df{kodimensi} dari suatu subruang $U$ dalam ruang vektor $V$ adalah dimensi dari~$V/U$. Jadi,
apabila $T$ suatu pemetaan linear, $\dim\coker T = \codim\ran T$.
\end{defn}

Dalam kategori $\cat{HIL}$ dari ruang Hilbert dan pemetaan linear kontinu, jangkauan suatu
morfisme tidak harus menjadi objek kategori tersebut. Secara khusus, jangkauan suatu operator
tidak harus tertutup.

\begin{exam}\label{004066} Operator ruang Hilbert
   \[ T\colon l_2 \sto l_2\colon (x_1,x_2,x_3,\dots) \mapsto
                                   (x_1, \tfrac12 x_2, \tfrac13 x_3, \dots) \]
bersifat injektif, swaadjoin, kompak, dan kontraktif,
 \index{operator!dengan jangkauan tak tertutup}%
tetapi jangkauannya, walaupun padat dalam $l_2$, tidak mencakup seluruh~$l_2$.
\end{exam}

Walaupun tidak langsung berkaitan dengan tujuan kita saat ini, inilah tempat yang tepat untuk mencatat
fakta bahwa jumlah dua subruang tertutup dari suatu ruang Hilbert tidak harus tertutup.

\begin{exam}\label{004071} Misalkan $T$ operator pada ruang Hilbert $l_2$ yang didefinisikan dalam
 \index{subruang!jumlah subruang tertutup tidak harus tertutup}%
contoh~\ref{004066}, $M = l_2 \oplus \{\vc 0\}$, dan $N$ graf dari~$T$. Maka $M$
dan $N$ keduanya merupakan subruang tertutup dari ruang Hilbert $l_2 \oplus l_2$, tetapi $M + N$
tidak tertutup.
\end{exam}

\begin{proof} Verifikasi contoh~\ref{004071} diperoleh dengan mudah dari hasil berikut
dan contoh~\ref{004066}. \ns
\end{proof}

\begin{prop} Misalkan $H$ suatu ruang Hilbert, $T \in \ofml B(H)$, $M = H \oplus \{\vc 0\}$, dan $N$
merupakan graf dari~$T$. Maka
 \begin{enumerate}
   \item[(a)] Himpunan $N$ merupakan subruang dari $H \oplus H$.
   \item[(b)] Operator $T$ injektif jika dan hanya jika $M \cap N = \{\,(\vc 0,\vc 0)\,\}$.
   \item[(c)] Jangkauan $T$ padat dalam $H$ jika dan hanya jika $M + N$ padat dalam $H \oplus H$.
   \item[(d)] Operator $T$ surjektif jika dan hanya jika $M + N = H \oplus H$.
 \end{enumerate}
\end{prop}

Kabar baik bagi teori operator Fredholm adalah bahwa operator dengan kokernel
berdimensi hingga secara otomatis mempunyai jangkauan tertutup.

\begin{prop}\label{00411} Jika suatu pemetaan linear terbatas $A\colon H \sto K$ antara ruang Hilbert
memiliki kokernel berdimensi hingga, maka $\ran A$ tertutup dalam~$K$.
\end{prop}

Dalam \emph{Alternatif Fredholm IIIb}~\ref{0040343}, kita mengamati bahwa syarat (1)
berlebih dan juga bahwa (2) berlaku untuk \emph{sebarang} operator ruang Banach dengan jangkauan tertutup.
Hal ini memungkinkan kita menyatakan kembali~\ref{0040343} secara lebih hemat.

\begin{prop}[Alternatif Fredholm IV]\label{00413}
 \index{alternatif Fredholm!versi IV}%
Jika $T$ suatu operator Riesz--Schauder pada ruang Banach, maka
 \begin{enumerate}
    \item $T$ mempunyai jangkauan tertutup dan
    \item  $\dim\ker T = \dim \ker T^* < \infty$.
 \end{enumerate}
\end{prop}













 
